\chapter{La matière venue des aliments}\label{ch:g6:matter-from-food}

Un veau de \qty{40}{kg} devient, en deux ans, une vache de
\qty{600}{kg}. Une demi-tonne de vache neuve --- chair, \omterm{def:g3:bones-and-muscles:skeleton}{os}, cuir,
cornes --- est apparue de quelque part. La vache a mangé de l'herbe~;
la matière nouvelle \emph{n'est pas de l'herbe}~; et pourtant, sans
l'herbe, rien de tout cela n'existerait. Ce que les \omterm{def:g1:animals-around-us:animal}{animaux} font de
leur nourriture est l'affaire de ce chapitre --- le côté
\omterm{def:g5:food-webs:roles}{consommateurs} du registre de \cref{ch:g6:plants-are-producers}.

\section{Les animaux bâtissent leur matière à partir de leur nourriture}

\begin{proposition}[Les consommateurs transforment la matière organique]\label{prop:g6:matter-from-food:transform}
Un \omterm{def:g1:animals-around-us:animal}{animal} bâtit sa propre \omterm{def:g6:plants-are-producers:matter}{matière organique} --- \omterm{def:g3:bones-and-muscles:muscle}{muscle}, \omterm{def:g3:bones-and-muscles:skeleton}{os}, graisse,
\omterm{def:g1:animals-around-us:covering}{plume}, \omterm{def:g1:animals-around-us:covering}{coquille} --- à partir de la \emph{\omterm{def:g6:plants-are-producers:matter}{matière organique} qu'il
mange}. La \omterm{def:g4:journey-of-food:digestion}{digestion} (\cref{def:g4:journey-of-food:digestion})
démonte les aliments en \omterm{def:g4:journey-of-food:digestion}{nutriments}~; le \omterm{prop:g4:journey-of-food:absorb}{sang} les livre
(\cref{prop:g4:heart-and-blood:carries})~; et les \omterm{def:g5:first-look-at-cells:cell}{cellules} du corps
(\cref{prop:g5:first-look-at-cells:growth}) les réassemblent en la
substance propre de l'\omterm{def:g1:animals-around-us:animal}{animal}. Les \omterm{def:g1:animals-around-us:animal}{animaux} sont des
\emph{transformateurs} de matière, jamais des fabricants à partir du
minéral~: ce métier appartient aux seuls \omterm{def:g5:food-webs:roles}{producteurs}.
\end{proposition}
\admitted

\begin{example}[De l'herbe en vache]\label{ex:g6:matter-from-food:cow}
Suis la demi-tonne~: la vache broute heure après heure
(\cref{prop:g2:what-animals-eat:daily} expliquait ces longues
heures)~; sa lourde \omterm{def:g4:journey-of-food:digestion}{digestion}
(\cref{exo:g4:journey-of-food:11}) démonte l'herbe en \omterm{def:g4:journey-of-food:digestion}{nutriments}~;
son \omterm{prop:g4:journey-of-food:absorb}{sang} les porte au pis, aux \omterm{def:g3:bones-and-muscles:muscle}{muscles} et au veau qui grandit en
elle~; ses \omterm{def:g5:first-look-at-cells:cell}{cellules} bâtissent de la matière de vache à partir de
matière d'herbe. Le pâturage se promène, très littéralement, sous
une forme nouvelle.
\end{example}

\begin{remark}[La matière change de mains, pas de nature]\label{rem:g6:matter-from-food:hands}
À chaque flèche d'une \omterm{def:g3:food-chains:chain}{chaîne alimentaire}
(\cref{def:g3:food-chains:chain}), la \omterm{def:g6:plants-are-producers:matter}{matière organique} change de
propriétaire et est rebâtie --- matière d'herbe en matière de vache,
matière de vache (le \omterm{def:g2:animals-grow-up:born}{lait}, la viande) en matière d'enfant.
L'\emph{organique} ne vient jamais de nulle part~: les \omterm{def:g5:food-webs:roles}{consommateurs}
se le passent et le remodèlent. Ce n'est qu'au départ vert de la
chaîne qu'il a été fabriqué neuf
(\cref{prop:g6:plants-are-producers:firstlink})~; ce n'est qu'au
bout, chez les \omterm{def:g5:food-webs:roles}{décomposeurs} du chapitre suivant, qu'il retourne au
minéral.
\end{remark}

\section{Grandir, c'est produire de la matière}

\begin{definition}[La production de matière d'un animal]\label{def:g6:matter-from-food:production}
La \emph{production de matière}\index{production de matière} d'un
\omterm{def:g1:animals-around-us:animal}{animal} est la \omterm{def:g6:plants-are-producers:matter}{matière organique} neuve qu'il bâtit~: la croissance
des jeunes (les kilogrammes gagnés par le veau), les renouvellements
annuels (le \omterm{def:g1:animals-around-us:covering}{pelage} mué et repoussé, les bois, les \omterm{def:g1:animals-around-us:covering}{plumes}), les
réparations (l'\omterm{def:g3:bones-and-muscles:skeleton}{os} ressoudé de
\cref{rem:g3:bones-and-muscles:alive}) et les productions (\omterm{def:g2:animals-grow-up:born}{lait},
\omterm{def:g2:animals-grow-up:born}{œufs}, laine). Tout cela est de la nourriture rebâtie.
\end{definition}

\begin{example}[Peser la production]\label{ex:g6:matter-from-food:weighing}
Les éleveurs pèsent la production chaque jour~: une vache laitière
donne quelque \qty{25}{kg} de \omterm{def:g2:animals-grow-up:born}{lait} par jour --- de la \omterm{def:g6:plants-are-producers:matter}{matière
organique} qu'elle doit rebâtir à partir de sa ration, jour après
jour~; l'\omterm{def:g2:animals-grow-up:born}{œuf} d'une poule représente quelque \qty{60}{g} de
provisions bien rangées (\cref{rem:g2:animals-grow-up:egg}),
assemblées dans la nuit à partir de son grain. Pas de ration, pas de
production~: le registre s'équilibre sans pitié.
\end{example}

\begin{proposition}[Toute la nourriture ne devient pas du corps]\label{prop:g6:matter-from-food:losses}
Un \omterm{def:g5:food-webs:roles}{consommateur} ne convertit jamais toute sa nourriture en matière
neuve~:
\begin{itemize}
  \item une part n'est jamais absorbée et sort en excréments
        (\cref{prop:g4:journey-of-food:absorb})~;
  \item une grande part est \emph{brûlée} --- combinée au
        dioxygène du souffle
        (\cref{prop:g4:breathing:exchange}) --- pour alimenter le
        mouvement et la chaleur~: de la matière dépensée comme
        carburant, qui repart en gaz expiré~;
  \item seul le reste devient croissance et productions.
\end{itemize}
De l'herbe qu'une vache mange, seule une fraction modeste finit en
vache neuve --- et c'est pourquoi chaque étage de la pyramide de
\cref{prop:g5:food-webs:pyramid} est tellement plus petit que celui
du dessous.
\end{proposition}
\admitted

\begin{example}[La pyramide, enfin expliquée]\label{ex:g6:matter-from-food:pyramid}
Les tonnes d'herbe d'un pré nourrissent une population de lapins~;
les kilogrammes des lapins nourrissent un renard ou deux. À chaque
flèche, la plus grande part de la matière est brûlée pour vivre ou
perdue, et seule une fraction est passée vers le haut sous forme de
corps. Les grands chasseurs sont rares
(\cref{prop:g5:food-webs:pyramid}) parce qu'ils vivent des
\emph{restes des restes} --- une arithmétique, pas de la malchance.
\end{example}

\section{Des menus, des dents, et un seul et même métier}

\begin{proposition}[Des menus différents, la même transformation]\label{prop:g6:matter-from-food:menus}
\omterm{def:g2:what-animals-eat:kinds}{Herbivore}, \omterm{def:g2:what-animals-eat:kinds}{carnivore}, \omterm{def:g2:what-animals-eat:kinds}{omnivore}
(\cref{def:g2:what-animals-eat:kinds}) désignent des
\emph{sources} différentes de \omterm{def:g6:plants-are-producers:matter}{matière organique}, pas des métiers
différents. La vache qui rebâtit de l'herbe et le renard qui rebâtit
du lapin mènent tous deux la transformation de
\cref{prop:g6:matter-from-food:transform} --- avec un équipement
assorti au menu
(\cref{prop:g2:what-animals-eat:match},
\cref{exo:g4:journey-of-food:11})~: l'herbe exige des broyeuses et
une \omterm{def:g4:journey-of-food:digestion}{digestion} immense, la viande récompense des dents qui saisissent
et un tube court.
\end{proposition}
\admitted

\begin{method}[Faire le bilan de matière d'un animal]\label{met:g6:matter-from-food:audit}
Pour n'importe quel \omterm{def:g1:animals-around-us:animal}{animal}, dresse son registre de matière~:
\begin{enumerate}
  \item \emph{entrées}~: quelle \omterm{def:g6:plants-are-producers:matter}{matière organique} mange-t-il, et à
        qui la prend-il~? (ses flèches entrantes du \omterm{def:g5:food-webs:web}{réseau
        alimentaire}, \cref{def:g5:food-webs:web})~;
  \item \emph{bâti}~: quelles croissances, quels renouvellements et
        quelles productions fabrique-t-il~?
  \item \emph{sorties}~: les excréments, et la matière brûlée pour
        le mouvement et la chaleur~;
  \item vérifie l'équilibre~: la production doit être le reste
        modeste --- si une affirmation la rend plus grande
        («~cet étang produit plus de matière de poisson que
        l'étang ne fait pousser de nourriture~»), l'affirmation est
        fausse.
\end{enumerate}
\end{method}

\begin{example}[Le bilan en pratique]\label{ex:g6:matter-from-food:audituse}
Fais le bilan du merle (\cref{ex:g2:what-animals-eat:winter})~:
entrées --- vers, \omterm{prop:g3:sorting-living-things:groups}{insectes}, baies~; bâti --- sa propre croissance,
les \omterm{def:g1:animals-around-us:covering}{plumes} neuves de chaque année, la ponte~; sorties --- des
fientes sur le mur, et la grande part brûlée qui a alimenté une
année de vol et de chant. Le poids de nourriture quotidien d'un
oiseau chanteur est énorme pour sa taille précisément parce que le
vol brûle une si grande part du registre.
\end{example}

\section{Exercices}

\begin{exercise}[$\star$]\label{exo:g6:matter-from-food:1}
Énonce le métier des \omterm{def:g5:food-webs:roles}{consommateurs}, et oppose-le à celui des
\omterm{def:g5:food-webs:roles}{producteurs}, en une phrase chacun.
\end{exercise}

\begin{exercise}[$\star$]\label{exo:g6:matter-from-food:2}
Énumère les étapes qui changent l'herbe en \omterm{def:g3:bones-and-muscles:muscle}{muscle} de vache, en
nommant le chapitre qui a traité chacune.
\end{exercise}

\begin{exercise}[$\star$]\label{exo:g6:matter-from-food:3}
Donne quatre formes de \omterm{def:g6:matter-from-food:production}{production de matière} chez un \omterm{def:g1:animals-around-us:animal}{animal}, avec un
exemple de chacune.
\end{exercise}

\begin{exercise}[$\star$]\label{exo:g6:matter-from-food:4}
Nomme les trois destinations de la nourriture d'un \omterm{def:g5:food-webs:roles}{consommateur},
d'après \cref{prop:g6:matter-from-food:losses}.
\end{exercise}

\begin{exercise}[$\star$]\label{exo:g6:matter-from-food:5}
Pourquoi une vache laitière a-t-elle besoin d'une ration chaque jour
pour continuer à donner du \omterm{def:g2:animals-grow-up:born}{lait}~?
\end{exercise}

\begin{exercise}[$\star$]\label{exo:g6:matter-from-food:6}
Un \omterm{def:g2:animals-grow-up:born}{œuf} contient \qty{60}{g} de \omterm{def:g6:plants-are-producers:matter}{matière organique}. Où a-t-elle été
fabriquée --- dans la poule, ou avant elle~? Démêle cela
soigneusement.
\end{exercise}

\begin{exercise}[$\star\star$]\label{exo:g6:matter-from-food:7}
Sers-toi de \cref{prop:g6:matter-from-food:losses} pour expliquer
pourquoi la pyramide alimentaire se rétrécit à chaque étage.
\end{exercise}

\begin{exercise}[$\star\star$]\label{exo:g6:matter-from-food:8}
Applique \cref{met:g6:matter-from-food:audit} à un chat de maison~:
entrées, bâti, sorties --- avec au moins deux éléments par ligne.
\end{exercise}

\begin{exercise}[$\star\star$]\label{exo:g6:matter-from-food:9}
«~Le renard et la vache font des choses différentes de leur
nourriture.~» Corrige la phrase avec
\cref{prop:g6:matter-from-food:menus} --- qu'est-ce qui diffère, et
qu'est-ce qui est identique~?
\end{exercise}

\begin{exercise}[$\star\star$]\label{exo:g6:matter-from-food:10}
Un enfant grandit de \qty{3}{kg} en un an alors qu'il mange
plusieurs centaines de kilogrammes de nourriture. Où est passé le
reste~? Réponds avec le registre.
\end{exercise}

\begin{exercise}[$\star\star$]\label{exo:g6:matter-from-food:11}
La \omterm{def:g1:animals-around-us:covering}{coquille} de l'escargot est une substance minérale, et pourtant
\cref{ex:g6:plants-are-producers:sorting} la comptait comme
l'ouvrage de l'escargot. Concilie les deux~: qu'a fait le corps de
l'escargot, et à partir de quelles fournitures, pour la bâtir~?
\end{exercise}

\begin{exercise}[$\star\star\star$]\label{exo:g6:matter-from-food:12}
Un hectare de terre nourrit plus de gens en pommes de terre qu'en
porc élevé sur ces mêmes pommes de terre. Explique avec
\cref{prop:g6:matter-from-food:losses} --- où passe la différence
quand les pommes de terre font le détour par le cochon~?
\end{exercise}

\section{Problème~: le registre de la ferme}

\begin{problem}[{Problème du week-end --- une petite ferme mise en
comptabilité, du champ à la table}]\label{pb:g6:matter-from-food:1}
La petite ferme de mamie~: un pâturage avec une vache, un poulailler
de dix poules nourries au grain, un carré de légumes, et la cuisine
familiale. Ce week-end, tu tiens son registre de matière.

\medskip\noindent\textbf{Partie I --- L'étage des \omterm{def:g5:food-webs:roles}{producteurs}.}
\begin{enumerate}
  \item Énumère les \omterm{def:g5:food-webs:roles}{producteurs} de la ferme. Quelles sont leurs
        entrées minérales, et leur source d'\omterm{prop:g3:balanced-diet:jobs}{énergie}~?
  \item Le pâturage produit environ \qty{10000}{kg} de matière
        d'herbe en une saison. Dans les \omterm{def:g1:plants-around-us:parts}{feuilles} de qui chaque
        gramme a-t-il été fabriqué~?
  \item Mamie épand le fumier du poulailler sur le carré de légumes.
        Quelle règle de
        \cref{prop:g5:people-and-nature:lasting} respecte-t-elle, et
        que réapprovisionne exactement le fumier~?
  \item Pourquoi ne peut-on pas simplement nourrir la vache de \omterm{prop:g4:what-plants-need:needs}{sels
        minéraux} et d'eau, en se passant du pâturage~?
\end{enumerate}

\medskip\noindent\textbf{Partie II --- L'étage des \omterm{def:g5:food-webs:roles}{consommateurs}.}
\begin{enumerate}[resume]
  \item La vache mange environ \qty{70}{kg} d'herbe par jour et
        donne à peu près \qty{20}{kg} de \omterm{def:g2:animals-grow-up:born}{lait}. Nomme les trois
        destinations des \qty{50}{kg} manquants.
  \item Les poules changent le grain en \omterm{def:g2:animals-grow-up:born}{œufs}. Écris le registre de
        la poule --- entrées, bâti, sorties --- et dis pourquoi une
        poule en mue (qui renouvelle toutes ses \omterm{def:g1:animals-around-us:covering}{plumes}) pond moins
        pendant un temps.
  \item Une buse emporte un poussin. Prolonge jusqu'à elle la
        \omterm{def:g3:food-chains:chain}{chaîne alimentaire} de la ferme, et dis à quels maillons la
        matière a été \emph{fabriquée} et à quels maillons elle n'a
        été que transformée.
  \item La famille mange des \omterm{def:g2:animals-grow-up:born}{œufs}, du \omterm{def:g2:animals-grow-up:born}{lait} et des légumes. Retrace
        la \omterm{def:g6:plants-are-producers:matter}{matière organique} d'un petit déjeuner --- pain, \omterm{def:g2:animals-grow-up:born}{œuf} à la
        coque, \omterm{def:g2:animals-grow-up:born}{lait} --- jusqu'à la \omterm{def:g1:plants-around-us:parts}{feuille} où chacun a d'abord été
        fabriqué.
\end{enumerate}

\medskip\noindent\textbf{Partie III --- L'arithmétique des
étages.}
\begin{enumerate}[resume]
  \item Mamie dit~: «~le pâturage nourrit une vache~; il ne
        nourrirait pas plus qu'une petite tanière de renards si des
        renards y mangeaient des lapins nourris d'herbe~». Justifie
        son estimation avec les pertes de chaque étage.
  \item Le carré de légumes nourrit la famille directement. En
        suivant la logique de \cref{exo:g6:matter-from-food:12},
        explique pourquoi le carré nourrit plus de gens par mètre
        carré que le trajet pâturage-vache-lait.
  \item En période de sécheresse, l'herbe cesse de pousser. Prédis,
        dans l'ordre, ce qui arrive au \omterm{def:g2:animals-grow-up:born}{lait} de la vache, à son
        poids, et à la ligne «~bâti~» du registre --- et pourquoi
        aucun surcroît de soleil ne peut aider tant que la pluie
        manque (\cref{pb:g6:plants-are-producers:1}, question 7,
        détient la clé).
  \item Referme le registre avec une phrase par étage~:
        \omterm{def:g5:food-webs:roles}{producteurs}, \omterm{def:g5:food-webs:roles}{consommateurs} --- et nomme ceux qui manquent
        encore à la comptabilité (le sujet du chapitre suivant).
\end{enumerate}
\end{problem}
